\documentclass[11pt]{article}
\usepackage[margin=1in]{geometry}
\usepackage{amsmath,amssymb}
\usepackage{booktabs}
\usepackage{hyperref}
\usepackage{siunitx}
\hypersetup{colorlinks=true,linkcolor=blue,urlcolor=blue}

\title{Independent Replication Report:\\
\emph{Theory of the Ferroelectric Phase in Organic Conductors:\\ Optics and Physics of Solitons}\\
\large S. Brazovskii, arXiv:cond-mat/0306006v2 (ECRYS-2002 proceedings)}
\author{REPLICATE-PROJECT / TEXTURES-100 --- independent reimplementation}
\date{2026-07-19}

\begin{document}
\maketitle

\begin{abstract}
This is an independent, from-scratch reimplementation of the analytic physics of
Brazovskii's ECRYS-2002 proceedings paper on the ferroelectric (FE) phase of the
quasi-1D organic conductors (TMTTF)$_2$X. The paper is a \emph{theory/proceedings}
paper: it presents closed-form results ``without derivations'' (its own words,
l.173) and contains \textbf{no numerical tables}. Accordingly the replication
scores agreement against \textbf{functional forms} and \textbf{order-of-magnitude}
claims, with every stated relation turned into a machine-checkable identity.
All ten analytic relations reimplemented pass (10/10), most to machine precision.

\medskip
\noindent\fbox{\parbox{0.96\linewidth}{\centering
\textbf{HEADLINE:} optical collective-mode edge relative to the photoconductivity
gap, $\ \omega_t/(2\Delta) = \pi\gamma/2 = \mathbf{0.393}$ at $\gamma=0.25$ ---
confirming the paper's claim that the optical edge $\omega_t\approx\pi\gamma\Delta$
lies \emph{well below} $2\Delta$ (l.22, l.153).\\[4pt]
\textbf{VERDICT: PARTIAL (mechanism-level REPLICATED).}\quad
Coverage $\mathbf{8/10}$ \quad Agreement $\mathbf{9/10}$.}}
\end{abstract}

\section{The paper's claims (as reimplemented)}
The paper describes a \emph{combined Mott--Hubbard} state of (TMTTF)$_2$X in which
the interference of a built-in bond nonequivalence ($U_b$) and a spontaneous site
nonequivalence ($U_s$, the charge disproportionation, CD) produces ferroelectricity.
The checkable analytic claims are:
\begin{enumerate}
\item Combined Mott--Hubbard Hamiltonian rewrites as a single cosine with an
      effective amplitude $U$ and phase shift $\alpha$; ground state at $\phi=\alpha$.
\item Three soliton species: $\pi$-solitons (holons, charge $e$); noninteger
      $\alpha$-solitons (FE domain walls, $q=-2\alpha/\pi$ or $1-2\alpha/\pi$);
      combined spin-charge solitons (short charge core, long spin tail).
\item Gap scaling $\Delta\sim U^{1/(2-2\gamma)}$; CD requires $\gamma<1/2$.
\item Optical edge is a \emph{collective mode} at $\omega_t\approx\pi\gamma\Delta<2\Delta$,
      not the two-particle gap $2\Delta$ (photoconductivity edge).
\item Mixed electron-phonon dielectric response $\varepsilon(\omega)$ [Eq.~2],
      exhibiting a Fano antiresonance at $\omega_0$, a combined resonance at
      $\omega_{0t}^2\approx\omega_0^2+\omega_t^2$, and an FE soft mode at
      $\omega_{fe}^2\approx(1-Z)/(\omega_0^{-2}+\omega_t^{-2})$.
\item Curie-law dielectric divergence $\varepsilon(T)=A|T/T_0-1|^{-1}$ with
      $A\sim(\omega_p^*/\omega_t)^2\sim10^3$.
\item Spin-Peierls spin gap $\Delta_\sigma\sim U_{ao}^{2/3}$; core/tail lengths
      $\xi_\rho\sim\hbar v_F/\Delta\ll\xi_\sigma\sim\hbar v_F/\Delta_\sigma$.
\end{enumerate}

\section{Governing equations reimplemented}
\paragraph{Combined Mott--Hubbard Hamiltonian.}
\begin{equation}
H_U = -U_s\cos 2\phi - U_b\sin 2\phi = -U\cos(2\phi-2\alpha),\quad
U=\sqrt{U_s^2+U_b^2},\ \ \tan 2\alpha = U_b/U_s .
\end{equation}
\paragraph{Gap and optical edge.}
$\Delta\sim U^{1/(2-2\gamma)}$, \ $\omega_t\approx\pi\gamma\Delta$, \ photoconductivity gap $=2\Delta$.
\paragraph{Dielectric response (Eq.~2).}
\begin{equation}
\frac{\varepsilon(\omega)}{\varepsilon_\infty}=1+
\frac{(\omega_p^*/\omega_t)^2\,(1-(\omega/\omega_0)^2)}
{(1-(\omega/\omega_0)^2)(1-(\omega/\omega_t)^2)-Z},\quad
Z=\left(\frac{\omega_{cr}}{\omega_t}\right)^{2-4\gamma}\le1 .
\end{equation}
The denominator is a quadratic in $x=\omega^2$:
$x^2-(\omega_0^2+\omega_t^2)x+(1-Z)\omega_0^2\omega_t^2=0$, whose two roots are the
combined resonance (upper) and the FE soft mode (lower).
\paragraph{Curie law.}
$\varepsilon(0)/\varepsilon_\infty = 1 + A/(1-Z)$, $A=(\omega_p^*/\omega_t)^2$; with
$1-Z\propto|T/T_0-1|$ this is the Curie $1/|t|$ divergence.

\section{Method}
Pure analytic / algebraic reimplementation in Python (NumPy 2.3.5), from the
equations transcribed above --- \textbf{not} author code (the paper ships none).
Each stated identity is evaluated numerically and asserted:
\begin{itemize}
\item Identities (H, soliton charges, Fano) checked to machine precision over grids.
\item The paper's $\approx$/$\sim$ formulas (resonance, soft mode) are the $Z\to1$
      (criticality, $T\to T_0$) leading-order limit of the \emph{exact} denominator
      roots; verified by scanning $Z=0.5\to0.9999$ and showing the relative error
      \emph{monotonically decreases} toward the limit.
\item Order-of-magnitude claims (Curie amplitude) checked by reproducing the exact
      Curie law and confirming $(\varepsilon-1)(1-Z)=A$ is constant.
\item Dimensionless ratios/exponents tested (the paper never fixes prefactors
      $C$, $\omega_p^*$, or the $T$-laws $\omega_t(T),\omega_{cr}(T)$).
\end{itemize}

\section{Results: this work vs.\ paper}
\begin{table}[h]
\centering
\begin{tabular}{@{}p{6.4cm}rrl@{}}
\toprule
Quantity / identity & This work & Paper & Status \\
\midrule
$H$ form identity $\max|{\rm LHS-RHS}|$ ($\phi$ grid) & $4.4\times10^{-16}$ & exact & \textbf{EXACT} \\
Ground-state min at $\phi=\alpha$ & true & $\phi=\alpha$ & \textbf{MATCH} \\
$\alpha$-soliton charge $-2\alpha/\pi$ ($U_s{=}0.7,U_b{=}0.4$) & $-0.1652$ & $-2\alpha/\pi$ & \textbf{EXACT} \\
$\alpha$-soliton charge $1-2\alpha/\pi$ & $0.8348$ & $1-2\alpha/\pi$ & \textbf{EXACT} \\
$\omega_t/(2\Delta)$ at $\gamma=0.25$ (=$\pi\gamma/2$) & $0.3927$ & ``$<2\Delta$'' & \textbf{MATCH} \\
$\gamma$ ceiling for $\omega_t<2\Delta$ ($=2/\pi$) & $0.6366$ & --- & consistent \\
Resonance $\omega_{0t}$ rel.\ err.\ at $Z{=}0.9999$ & $1.2\times10^{-5}$ & $\approx$ limit & \textbf{CONVERGES} \\
Soft mode $\omega_{fe}$ rel.\ err.\ at $Z{=}0.9999$ & $1.2\times10^{-5}$ & $\approx$ limit & \textbf{CONVERGES} \\
Fano antiresonance $\varepsilon/\varepsilon_\infty(\omega_0)$ & $1.0$ & $1$ (antires.) & \textbf{EXACT} \\
Curie amplitude $A=(\omega_p^*/\omega_t)^2$ & $900$ & $\sim10^3$ & \textbf{ORDER MATCH} \\
Curie product $(\varepsilon-1)(1-Z)$ constancy (std/mean) & $5.5\times10^{-16}$ & $1/|t|$ law & \textbf{EXACT} \\
Soft mode $\propto\sqrt{t}$ (corr.\ coeff.) & $>0.999$ & $\sim\varepsilon^{-1/2}$ & \textbf{MATCH} \\
Spin tail / core ratio $\xi_\sigma/\xi_\rho$ ($U_{ao}{=}0.05$) & $7.37$ & $\gg1$ & \textbf{MATCH} \\
\bottomrule
\end{tabular}
\caption{All 10 verdict checks pass. Machine-precision identities plus
monotone convergence of the two $\approx$-formulas to their criticality limit.}
\end{table}

\paragraph{Convergence of the $\approx$-formulas (the key methodological point).}
Testing $\omega_{0t}^2\approx\omega_0^2+\omega_t^2$ at an \emph{arbitrary} $Z$
(e.g.\ $Z=0.5$) gives $\sim7\%$ ``error'' and a false miss. It is a limit formula:
scanning $Z=0.5,0.9,0.99,0.999,0.9999$ gives resonance relative errors
$7.0\%\to1.2\%\to0.12\%\to0.012\%\to0.0012\%$ (and identically for the soft mode).
Monotone convergence toward $T_0$ \emph{is} the verification.

\section{Comparison \& verdict}
\textbf{PARTIAL (mechanism-level REPLICATED); coverage 8/10, agreement 9/10.}
Every algebraic identity in Sections I--II of the paper reproduces exactly
(4$\times10^{-16}$ on the Hamiltonian form, Fano $=1$, Curie $1/|t|$ exact to
$5\times10^{-16}$). The two $\approx$ formulas are confirmed as the leading
criticality limit of exact objects. The order-of-magnitude Curie amplitude
($A=900\sim10^3$) matches the paper's own estimate; the experimental $\sim10^4$
is the number the paper itself only approximately matches, not a target this code
must hit. Agreement is 9 (not 10) only because there are no numerical tables to
fit against --- the score is necessarily on forms and orders of magnitude.

\section{Honest gaps and scope}
\begin{enumerate}
\item \textbf{Proceedings paper, no tables.} Results are given ``without derivations''
      (l.173); agreement is on functional forms + order-of-magnitude, not tabulated data.
      \emph{Expected, not a shortfall.}
\item \textbf{Free normalizations.} The prefactor $C$ in $\Delta\sim U^{1/(2-2\gamma)}$,
      the plasma frequency $\omega_p^*$, and the temperature laws
      $\omega_t(T),\omega_{cr}(T)$ are not fixed by the paper; representative values
      were chosen. Only dimensionless ratios and scaling exponents are tested.
\item \textbf{Quantum breather spectrum} between $\omega_t$ and $2\Delta$
      (soliton-antisoliton bound states) is \emph{mentioned but given no closed form};
      not reimplemented. \emph{Coverage-capping, scoped out.}
\item \textbf{Microscopic mass renormalization} ($U\to U^*$, $\gamma$ from
      $v_\rho/v_F$, ripplon analogy) given only as scaling relations; checked
      structurally, not against data.
\end{enumerate}
Items 3--4 are exactly why coverage caps at 8/10 for this proceedings paper.

\section{Reproducibility}
Interpreter: \texttt{/home/stevens/comfyui-env/bin/python} (NumPy 2.3.5, SciPy 1.17.0).
The solver is \texttt{report/evidence/brazovskii2003\_replication.py}; the live re-run
reproduces \texttt{report/evidence/brazovskii2003\_result.json} (and
\texttt{brazovskii2003\_output.json}) to the digits quoted above (10/10 checks pass).
\begin{verbatim}
cd textures-polar-brazovskii2003/
/home/stevens/comfyui-env/bin/python report/evidence/brazovskii2003_replication.py
\end{verbatim}
No LaTeX engine is installed on the packaging host; this \texttt{.tex} is delivered as
source and compiles off-host with a standard \texttt{pdflatex} run.

\end{document}
